
\chapter{Conclusion générale et perspectives}

\section{Bilan du projet}

Au terme de ce projet, nous avons réalisé une application web de gestion des consultations médicales en utilisant les technologies Java EE et les outils associés au développement d'applications web.

Ce travail nous a permis de mettre en pratique les connaissances acquises durant notre formation, notamment dans le développement côté serveur avec Java, l'utilisation des Servlets et des JSP, ainsi que la gestion des données à l'aide de MySQL et Hibernate.

L'application a été développée en adoptant une architecture MVC, avec une séparation des responsabilités entre les différentes couches DAO, Service, Servlet et JSP. Cette organisation a permis de structurer le code et de faciliter la communication entre les différentes composantes du système.

Par ailleurs, ce projet nous a permis de découvrir les différentes étapes de réalisation d'une application web, depuis l'analyse des besoins et la conception jusqu'à l'implémentation des fonctionnalités.

Ainsi, ce projet constitue une expérience pratique qui nous a permis de mieux comprendre le fonctionnement et l'organisation d'une application web basée sur Java.

\section{Difficultés rencontrées}

Durant la réalisation de ce projet, nous avons rencontré plusieurs difficultés techniques liées à la mise en place et à l'intégration des différentes technologies utilisées.

Parmi les principales difficultés rencontrées, nous pouvons citer :

\begin{itemize}
	\item La configuration de l'environnement de développement et l'intégration des différentes dépendances nécessaires au fonctionnement du projet.
	
	\item La configuration de la connexion entre l'application Java et la base de données MySQL à travers Hibernate et JPA.
	
	\item La compréhension et la mise en place de la communication entre les différentes couches de l'architecture MVC.
	
	
	\item La résolution des erreurs rencontrées durant le développement et l'intégration des différentes composantes de l'application.
\end{itemize}

Ces difficultés nous ont amenés à approfondir notre compréhension des technologies utilisées et à améliorer progressivement notre manière de développer et d'organiser une application web.

\section{Compétences acquises}

La réalisation de ce projet nous a permis de développer plusieurs compétences techniques et méthodologiques.

Sur le plan technique, nous avons renforcé nos connaissances en programmation Java et dans le développement d'applications web avec Jakarta EE. Nous avons également acquis une expérience pratique dans l'utilisation des Servlets, des JSP, de JDBC et d'Hibernate pour assurer la communication entre les différentes composantes de l'application.

Nous avons également approfondi nos connaissances en conception et en gestion des bases de données relationnelles avec MySQL, ainsi que dans l'utilisation de l'architecture MVC et de la séparation en couches DAO et Service.

En complément, ce projet nous a permis de nous familiariser davantage avec les outils de développement tels qu'IntelliJ IDEA, Apache Maven, Apache Tomcat, DBeaver et GitHub.

Enfin, le travail en équipe nous a permis de développer notre capacité à collaborer, à partager le code source et à coordonner les différentes tâches nécessaires à la réalisation d'un projet informatique.

\section{Améliorations possibles}

Bien que l'application réalisée réponde aux objectifs définis dans le cadre de ce projet, plusieurs améliorations pourraient être envisagées afin d'enrichir ses fonctionnalités et d'améliorer son utilisation.

Parmi les perspectives d'évolution possibles, nous pouvons citer :

\begin{itemize}
	\item L'ajout d'un système de notifications pour informer les utilisateurs des nouvelles consultations et des modifications apportées aux rendez-vous.
	
	\item La mise en place d'un calendrier interactif permettant de visualiser et de gérer plus facilement les consultations.
	
	\item L'amélioration de la sécurité à travers un système de gestion des rôles et des autorisations plus avancé.
	
	\item L'ajout de fonctionnalités de statistiques et de tableaux de bord pour faciliter le suivi de l'activité médicale.
	
	\item L'amélioration de l'interface utilisateur afin de rendre l'application plus ergonomique et accessible sur différents appareils.
	
	\item La mise en place de tests automatisés afin de vérifier le bon fonctionnement des différentes fonctionnalités.
\end{itemize}

Ces améliorations pourraient contribuer à rendre l'application plus complète, plus sécurisée et mieux adaptée aux besoins de ses futurs utilisateurs.